\documentclass[11pt]{article}
\usepackage[utf8]{inputenc}
\usepackage[letterpaper,top=0cm, margin=0.85in]{geometry}

\usepackage{textcmds} %more symbols
\usepackage{fontspec} %more fonts


%for math
\usepackage{amsmath, amstext, amssymb, amsfonts} %standard
\usepackage{youngtab} % makes squares for math diagrams
\usepackage{microtype} %% <-- added
%-----------------------------------------------------------           

%\usepackage{sectsty}
%for lists and numbers
\usepackage{enumitem}
%-----------------------------------------------------------

% Doc setting
\usepackage[english]{babel} % Replace `english' with e.g. `spanish' to change the document language
\usepackage{setspace} %to set spacing bw words and lines
\usepackage{changepage}
% \setlength\parindent{0pt}

%footer
\usepackage{fancyhdr}
\usepackage{lastpage}

% \fancyhf{0pt} % sets both header and footer to nothing
% \renewcommand{\footrulewidth}{0pt}
% \renewcommand{\headrulewidth}{0pt} %remove headerline

% \fancyfoot[RE,RO]{\thepage}
\fancyfoot[C]{MATH 201 | \emph{Shubhro Gupta}}
\renewcommand{\headrulewidth}{0pt}% Default \headrulewidth is 0.4pt
\renewcommand{\footrulewidth}{0.4pt}% Default \footrulewidth is 0pt
\fancyfoot[L]{Assignment 2}
\fancyfoot[R]{\thepage}
\pagestyle{fancy}
%-----------------------------------------------------------

%for pictures and graphs
\usepackage{graphicx} %add image
\usepackage{adjustbox}

\usepackage{pgfplots} %for graphing plotting
\pgfplotsset{compat=1.18, width=10cm}
%-----------------------------------------------------------

%for code
\usepackage{verbatim}
\usepackage{listings}
\usepackage{fancyvrb} %for coding blocks
%\usepackage{algorithm}
%\usepackage{algpseudocode} %for pseudocode
%\usepackage{algorithm, algpseudocode}

%\usepackage{lstfiracode} %firacode
\usepackage[framemethod=tikz]{mdframed} %adding background to lstlisting
\usepackage[ruled,vlined,boxed]{algorithm2e} %for pseudocode lines



%for colors and links
\usepackage[colorlinks = true,
            linkcolor = blue,
            urlcolor  = blue,
            citecolor = blue,
            anchorcolor = blue]{hyperref}
\usepackage[many]{tcolorbox}  % for colored boxes
\usepackage{color} % to get colors
\usepackage{xcolor} %more colors options and flexibility
\usepackage{transparent}


%-----------------------------------------------------------------------------
%custom commands

%code
\newcommand{\problem
}[2]{
\begin{mdframed}
    \textsc{question} \textbf{#1} \hfill #2
\end{mdframed}
}
\newcommand{\codecap}[2]{{\vspace{4pt}{\emph{#1}}} \hfill \href{#2}{Link to the code\ }\vspace{25pt}}
\newcommand{\code}[1]{{\texttt{#1}}}

%math
\newcommand{\bigo}[1]{$O(#1)$ }
\newcommand{\thetan}[1]{$\theta(#1)$}
% \newcommand{\vector}[1]{$\overrightarrow{#1}$}}o
\newcommand{\contradiction}{%
\begin{tikzpicture}[rotate=45,x=0.5ex,y=0.5ex]
\draw[line width=.2ex] (0,2) -- (3,2) (0,1) -- (3,1) (1,3) -- (1,0) (2,3) -- (2,0);
\end{tikzpicture}
}

\newcommand{\vecset}[2]{\{ {#1}_1, {#1}_2, {#1}_3,  \dots,  {#1}_{#2}\}}

%display
\newcommand{\link}[3][blue]{\href{#2}{\color{#1}{#3}}}%
\newcommand{\inlink}[1]{\underline{\emph{\link[black]{#1}{#1}}}}


%header
\newcommand{\heading}[5]{
\hrule ~\\~\\
\begin{large}
\noindent\emph{#1}\smallskip ~\\
Professor #3 \hfill Week #2 \smallskip ~\\
\textbf{Shubhro Gupta} \hfill Due #4 ~\\
\end{large} \medskip ~\\
{\emph{Collaborators: #5}}~\\
\hrule
\vspace{50pt}
~\\
}

% \newcommand\dunderline[3][-1pt]{{%
%   \sbox0{#3}%
%   \ooalign{\copy0\cr\rule[\dimexpr#1-#2\relax]{\wd0}{#2}}}}


%-----------------------------------------------------------------------------
%title
\usepackage{algpseudocode}
\begin{document}

\heading{Real Analysis}{1}{Rishi Vyas}{Monday, 18 August}{Khushi Tyagi}

\problem{A}{2 Points}
\textbf{To Show.} The set $(- \infty, a)$ is open, i.e., $B(x, \delta) \subseteq (- \infty, a)$ for some $ x \in (- \infty, a), \delta > 0$.\\
\textbf{Proof.} Let $x \in (- \infty, a)$, then $- \infty < x < a$.\\
Let $\delta = min(x + \infty, a - x)$.\\
This means $\delta \leq x + \infty$ and $\delta \leq a - x$.\\
We have to show that $B(x - \delta, x + \delta) \subseteq (- \infty, a)$.\\
From above $x - \delta \leq x  - (x + \infty) = - \infty$ and $x + \delta \leq x + (a - x) = a$. (By subtracting $x$ from both sides).\\
Thus, $B(x - \delta, x + \delta) \subseteq (- \infty, a)$.\\
Since for any $x \in (- \infty, a)$, we can find a $\delta > 0$ such that $B(x - \delta, x + \delta) \subseteq (- \infty, a)$, the set $(- \infty, a)$ is open. \hfill $\square$\\
\vspace{10pt}

\noindent \textbf{To Show.} The set $(\infty, a]$ is not open, i.e., $B(x, \delta) \not\subseteq (\infty, a]$ for any $ x \in (\infty, a], \delta > 0$.\\
\textbf{Proof.} Let $x \in (\infty, a]$, then $x \leq a < \infty$.\\
Assume $x = a$, and $B(x, \delta) \subseteq (\infty, a] \implies B(a - \delta, a + \delta) \subseteq (\infty, a]$.\\
This means $a + \delta \leq a \implies \delta \leq 0$.\\
But $\delta > 0$, hence $B(x, \delta) \not\subseteq (\infty, a]$.\\
Thus, for any $x \in (\infty, a]$, we cannot find a $\delta > 0$ such that $B(x, \delta) \subseteq (\infty, a]$. Hence, the set $(\infty, a]$ is not open. \hfill $\square$\\




\problem{B}{2 Points}
\textbf{To Determine. } All $x \in \mathbb{R}$ that satisfies $|x + 3| + |x - 2| = 9$.\\
\textbf{Solution.} \\
\underline{Case 1:} $x \geq 2$.\\
Then $|x + 3| = x + 3$ and $|x - 2| = x - 2$.\\
Thus, $x + 3 + x - 2 = 9 \implies 2x + 1 = 9 \implies x = 4$.\\
\\
\noindent \underline{Case 2:} $-3 \leq x < 2$.\\
Then $|x + 3| = x + 3$ and $|x - 2| = -x + 2$.\\
Thus, $x + 3 - x + 2 = 9 \implies 5 = 9$ which is not possible.\\
\\
Thus, there are no solutions for $-3 \leq x < 2$. \hfill $\square$\\




\problem{C}{2 Points}
\textbf{To Show. } $\mathbb{Z}^c$ is open in $\mathbb{R}$.\\
\textbf{Proof.}\\
$\mathbb{Z}^c = \underset{n \in \mathbb{Z}}{\bigcup} (n, n + 1)$.\\
Let $x \in \mathbb{Z}^c$, then $x \in (n, n + 1)$ for some $n \in \mathbb{Z}$.\\
Let $\delta = min(x - n, n + 1 - x)$.\\
Then $x - \delta \leq x - (x - n) = n$ and $x + \delta \leq x + (n + 1 - x) = n + 1$.\\
$B(x, \delta) = (x - \delta, x + \delta)$.
Consider $x - \delta$. If $x - \delta = n$, then $x - \delta \in (n, n + 1)$.\\
If $x - \delta < n$, then $x - \delta \in (n - 1, n)$.\\
Similarly, $x + \delta \in (n, n + 1)$ or $(n, n + 1)$.\\

\noindent Therefore $B(x, \delta) = (x - \delta, x + \delta) \subseteq (n, n + 1) \subseteq \mathbb{Z}^c$.\\


\problem{E}{2 Points}
\textbf{To Show.} sup$(S) \not\in S$ for any non-empty, bounded above set $S \subseteq \mathbb{R}$.\\
\textbf{Proof.} Assume sup$(S) \in S$, and sup$(S) = x$.\\
Since $S$ is open, $\exists \delta > 0$ such that $B(x, \delta) = (x - \delta, x + \delta) \subseteq S$.\\
Since $x$ is the supremum of $S$, $\forall s \in S, s \leq x$.
\\ $x < x + \delta$, since $\delta > 0$.\\
Also, $x + \delta \in S$, since $x + \delta \in B(x, \delta) \subseteq S$.\\
But $x + \delta > x$, which contradicts the fact that $x$ is the supremum of $S$.\\
Thus, sup$(S) \not\in S$ for any non-empty, bounded above set $S \subseteq \mathbb{R}$. \hfill $\square$\\






\problem{F}{4 Points}
\textbf{To Find. } Closure and limit of\\
(1) (0, 3] $\subseteq \mathbb{R}$\\
\textbf{Solution.} \\
\underline{Closure} $\overline{(0, 3]} = (0, 3] \cup \{0\}$.\\
We have to show that $B(x, \delta) \cap (0, 3] \neq \emptyset$ for any $x \in \overline{(0, 3]}$ and $\delta > 0$.\\
Let $y = min \frac{0 + \delta}{2}, 1.5$.\\
Since $0 < y < 1.5$, $y \in (0, 3]$.\\
Taking $x = 0$, we have $B(0, \delta) \cap (0, 3]$, we get $y \in B(0, \delta) \cap (0, 3]$.\\
Thus, $\overline{(0, 3]} = (0, 3] \cup \{0\}$.\\

\noindent \underline{Limit}
$\lim_{x \to 3} (0, 3] = 3$.\\
We have to show that $\forall \epsilon > 0, \exists \delta > 0$ such that $0 < |x - 3| < \delta \implies |f(x) - 3| < \epsilon$.\\
Let $\delta = \epsilon$.\\
Then $0 < |x - 3| < \delta \implies 0 < x - 3 < \delta \implies 0 < x - 3 < \epsilon \implies |f(x) - 3| < \epsilon$.\\
Thus, $\lim_{x \to 3} (0, 3] = 3$.\\

(2) $\mathbb{Z}$ Closure is $\mathbb{Z}$.\\
\textbf{Solution.} \\

(3) $\mathbb{Q} \cap (0, 1)$ Closure is $[0, 1]$.\\



\problem{G}{4 Points}
\textbf{Construct Example. } Of an infinite susbet of $\mathbb{R}$ which is not an interval, and is neither closed or open.\\
\textbf{Solution.} Consider the set $S = \{1\} \cup (2, 3)$.\\
$S$ is not an interval since it is not connected.\\
$S$ is not closed since it does not contain its limit points.\\
$S$ is not open since it does not contain all its boundary points.\\
\\
\textbf{Construct Example. } Of a set $S \subseteq \mathbb{R}$ which is bounded and has exactly 2 limit points.\\
\textbf{Solution.} Consider the set S =
$S$ is bounded since $0 < S < 3$.\\
The limit points of $S$ are 1 and 2, since




\problem{H}{3 Points}
Let $S \subseteq \mathbb{R}$ and $x \in \mathbb{R}$.\\
\textbf{To Show. } $x = \text{inf}(S) \iff $ The element $x$ is a lower bound and given $n \in \mathbb{N}, \exists y \in S$ such that $y < x + \frac{1}{n}$.\\
\textbf{Proof.} \\
\underline{Case 1} $x = \text{inf}(S)$.\\
Then $x$ is a lower bound of $S$ by definition.\\
If $x$ is an infimum, then $\forall 1 > \epsilon > 0, \exists y \in S$ such that $x \leq y < x + \epsilon$.\\
Since $\epsilon > 0$, we can write $\epsilon = \frac{1}{n}$ for some $n \in \mathbb{N}$.\\
Thus, $\forall n \in \mathbb{N}, \exists y \in S$ such that $y < x + \frac{1}{n}$.\\

\noindent \underline{Case 2} The element $x$ is a lower bound and given $n \in \mathbb{N}, \exists y \in S$ such that $y < x + \frac{1}{n}$.\\
Let $t$ be any lower bound of $S$.\\
Then $t \leq x$ since $x$ is the greatest lower bound.\\
We have to show that $x \leq t$. Suppose $x > t$.\\
Then $x - t > 0$.\\
Let $\epsilon = x - t$.\\
Then $\epsilon > 0$.\\
Since $\epsilon > 0$, we can write $\epsilon = \frac{1}{n}$ for some $n \in \mathbb{N}$.\\
Thus, $\exists y \in S$ such that $y < x + \frac{1}{n} = t + \epsilon$. $y - \epsilon < t$. But $y - \epsilon < x$.\\
This contradicts the fact that $x$ is the greatest lower bound.\\


\problem{I}{7 Points}
(1) \textbf{To Draw. } A picture of the set $S :=  \{ \frac{1}{n} - \frac{1}{m} : n, m \in \mathbb{N} \}$.\\
\textbf{Solution.} \\
(1) The set $S$ is the set of all possible differences between reciprocals of natural numbers.\\
(2) \textbf{To Show. } The supremum and infimum of $S$ are 1 and 0 respectively.\\
Let $x = \frac{1}{n} - \frac{1}{m} \in S$.\\
Then $0 < x < 1$.\\
Thus, $0$ is a lower bound of $S$.\\
Let $y \in S$, then $y = \frac{1}{n} - \frac{1}{m}$.\\
Then $y < 1$ for all $n, m \in \mathbb{N}$.\\
Thus, $1$ is an upper bound of $S$.\\

(3) \textbf{To Show. } The set $S$ is not open or closed.
\textbf{Proof.} \\
\underline{Open} Let $x \in S$, then $x = \frac{1}{n} - \frac{1}{m}$.\\
Let $\delta = min(\frac{1}{n}, \frac{1}{m})$.\\
Then $B(x, \delta) = (\frac{1}{n} - \delta, \frac{1}{m} + \delta)$.\\
Since $\delta < \frac{1}{n}$ and $\delta < \frac{1}{m}$, $B(x, \delta) \not\subseteq S$.\\
Thus, $S$ is not open.\\

\noindent \underline{Closed} Let $x \in S$, then $x = \frac{1}{n} - \frac{1}{m}$.\\
Let $\delta = min(\frac{1}{n}, \frac{1}{m})$.\\
Then $B(x, \delta) = (\frac{1}{n} - \delta, \frac{1}{m} + \delta)$.\\
Since $\delta < \frac{1}{n}$ and $\delta < \frac{1}{m}$, $B(x, \delta) \not\subseteq S$.\\
Thus, $S$ is not closed.\\

(4) \textbf{To Calculate. } The Closure.
\textbf{Solution.} \\
The closure of $S$ is $\overline{S} = S \cup \{0, 1\}$.\\








\end{document}
